\lysh{36896}{2}

\begin{alist}
\item Για $\lambda = 0$, η εξίσωση γίνεται: $-1 \cdot x = -1$.
Για $\lambda = 1$, η εξίσωση γίνεται: $0 \cdot x = 0$.
Για $\lambda = 3$, η εξίσωση γίνεται: $2x = 8$.

\item
\begin{rlist}
\item Η (1) έχει μια και μοναδική λύση αν και μόνο αν $\lambda - 1 \ne 0 \Leftrightarrow \lambda \ne 1$.
\item Για $\lambda \ne 1$, η μοναδική λύση της εξίσωσης είναι:
\[x = \dfrac{\lambda^2 - 1}{\lambda - 1} = \dfrac{(\lambda - 1)(\lambda + 1)}{\lambda - 1} = \lambda + 1.\]

Η λύση ισούται με $4$, οπότε $\lambda + 1 = 4 \Leftrightarrow \lambda = 3$.
Άρα για $\lambda = 3$ η λύση της εξίσωσης ισούται με $4$.
\end{rlist}
\end{alist}

